% =============================================================
%  Kapitel 3 – Verteilungsarchitektur (Deployment)
% =============================================================
\section{Verteilungsarchitektur (Deployment)}
\label{sec:deployment}

Das Deployment-Diagramm (Abbildung~\ref{fig:deployment}) ordnet Softwarekomponenten
den physischen Hardwareeinheiten zu und macht die Kommunikationskanäle sichtbar.

\begin{figure}[H]
  \centering
  \includegraphics[width=0.95\textwidth]{03_Deployment}
  \caption{Deployment-Diagramm – Hardware/Software-Mapping}
  \label{fig:deployment}
\end{figure}

\subsection{Komponentenverteilung}

\begin{center}
\renewcommand{\arraystretch}{1.3}
\begin{tabular}{>{\bfseries}p{4cm} p{9cm}}
  \toprule
  Hardware-Knoten & Softwarekomponenten \\
  \midrule
  Jetson Nano
    & \texttt{AutonomyController}, \texttt{TTCObserver},
      \texttt{SensorFusionAdapter} \\
  Arduino Mega Master
    & \texttt{MasterLanLink}, \texttt{MasterCanGateway},
      \texttt{HmiService} \\
  Arduino Mega Slave
    & \texttt{EnvironmentControl}, \texttt{EmergencyStop},
      \texttt{DriveControlUnit}, \texttt{Logger} \\
  \bottomrule
\end{tabular}
\end{center}

\subsection{Kommunikationskanäle}

\begin{itemize}
  \item \textbf{LAN} (Ethernet): Jetson Nano $\leftrightarrow$ Arduino Mega Master. \\
    Überträgt Autonomie-Kommandos und Zonen-Updates vom Jetson zum Master.
  \item \textbf{CAN-Bus}: Arduino Mega Master $\leftrightarrow$ Arduino Mega Slave. \\
    Überträgt CAN-Pakete mit Zonenstatus, Emergency-Ereignissen und Steuerbefehlen.
    \assumption{CAN-Datenklasse kapselt Paketstruktur (ID, Payload, Timestamp).}
\end{itemize}

\subsection{Rationale der Verteilung}

Die Platzierung des Emergency Stop direkt auf dem Arduino Mega Slave
neben der \texttt{DriveControlUnit} verfolgt folgende Ziele:

\begin{enumerate}
  \item \textbf{Minimale Latenz:} Kein zusätzlicher Kommunikationsweg zwischen
        Sicherheitslogik und Motorsteuerung.
  \item \textbf{Isolation:} Ausfall des Jetson oder des Masters hat keinen direkten
        Einfluss auf die Fähigkeit, einen Stopp einzuleiten.
  \item \textbf{Determinismus:} Der Slave führt eine einfache Echtzeit-Logik aus;
        kein OS-Overhead wie auf dem Jetson.
\end{enumerate}

\todobox{Latenzzahlen (typische CAN-Zykluszeit, LAN-Latenz) ergänzen,
sobald Messungen vorliegen.}
